% Generated draft: requires mathematical and source review.
\documentclass[11pt]{article}
\usepackage[margin=25mm]{geometry}
\usepackage{fontspec}
\setmainfont{Latin Modern Roman}
\usepackage{amsmath,amssymb,mathtools,unicode-math}
\setmathfont{Latin Modern Math}
\usepackage{enumitem,needspace,xurl,hyperref}
\setlength{\parindent}{0pt}
\setlength{\parskip}{5pt}
\setlength{\emergencystretch}{3em}
\newcommand{\N}{\mathbb{N}}
\newcommand{\Z}{\mathbb{Z}}
\newcommand{\Q}{\mathbb{Q}}
\newcommand{\R}{\mathbb{R}}
\newcommand{\C}{\mathbb{C}}
\newcommand{\F}{\mathbb{F}}
\newcommand{\A}{\mathbb{A}}
\newcommand{\PP}{\mathbb{P}}
\newcommand{\E}{\mathbb{E}}
\newcommand{\eps}{\varepsilon}
\newcommand{\dd}{\,\mathrm{d}}
\begin{document}
\section*{UnsolvedMath ID: 30000484}
\textbf{Source problem code:} catalogue entry
\subsection{Definitions and mathematical statement}
Does the bijection between noncrossing partitions and finitely generated wide subcategories extend to wild representation type?

A wide subcategory of an abelian module category is ordinarily a full subcategory closed under kernels, cokernels, and extensions, with the inherited exact abelian structure. The source does not identify the ambient algebra or quiver, the module category, or its notion of finite generation for a wide subcategory. Generation as the smallest wide subcategory containing finitely many objects and existence of a specified finite generator are not silently identified.

The ``noncrossing partitions'' likewise need an associated combinatorial object. In representation-theoretic correspondences they may be generalized objects attached to a Coxeter system rather than ordinary set partitions, but no Coxeter system, Coxeter element, order interval, or map is supplied here. The class in wild type must be defined before its cardinality or correspondence with subcategories can be asserted.

The request is to extend the source's specific correspondence from its known finite- and tame-type settings to the intended wild-type setting. This is not a universal assertion for every wild algebra: the extraction does not state whether hereditary algebras, acyclic quivers, or another class is intended, or specify the field or representation-type conventions. Nor does it state whether the extension should preserve an order, recover a given construction, or merely give a bijection. No arbitrary set-theoretic bijection is substituted for the missing report correspondence.

Thus the source question is retained with explicit missing definitions of its two sides and map. The dated assessment calls the wild-type extension conjectural; it does not supply the established finite/tame hypotheses or a complete wild formulation.
\subsection{Short English statement}
The question asks whether a representation-theoretic noncrossing-partition correspondence continues to classify finitely generated wide subcategories in wild type. The algebra class, partition model, finite-generation condition, and actual correspondence are not supplied.
\subsection{Sources}
\begin{itemize}
\item Local UnsolvedMath registry, permanent ID 30000484, code catalogue entry, ``Noncrossing Partitions in Wild Representation Type.'' Proposed year 2006; proposer unspecified.
\item \emph{Interactions between Algebraic Geometry and Noncommutative Algebra} (2007), Oberwolfach Reports, pp.\ 1319--1384, DOI 10.4171/owr/2006/23; original Report 23/2006 \url{https://doi.org/10.4171/owr/2006/23}. The displayed report year 2007 differs from the DOI and proposed year 2006.
\item Supplied reference labeled ``recent work on categories of noncrossing partitions'': \url{https://doi.org/10.1007/s10485-025-09838-8}. Its precise relevance to the representation-theoretic correspondence is not supplied.
\item Statement status exact; the supplied verification says official-report content was checked and notation, context, and question form normalized without changing the source claim.
\item Literature triage checked 2026-08-21: open. It says the finite/tame-to-wild extension is explicitly conjectural and not established by the later literature located in the supplied search. Created 2026-08-25; updated 2026-09-16. The underlying hypotheses and current status require independent verification.
\item Registry curation by Ulam AI / UnsolvedMath Contributors, CC BY 4.0; underlying report and literature provenance remain as supplied. Unreviewed statement-only draft.
\end{itemize}
\end{document}
